%=============================================================================
% Chapter 1 Solutions
%=============================================================================
\newpage
\section*{Chapter 1 Solutions}
\addcontentsline{toc}{section}{Chapter 1 Solutions}

%-----------------------------------------------------------------------------
\subsection*{1.1 --- Types of Computers}
\hypertarget{sol-1.1}{}
\markright{1.1}
\addcontentsline{toc}{subsection}{1.1 --- Types of Computers}

Three major types:

\begin{enumerate}
  \item \textbf{Personal computers (PCs)} --- general-purpose computers
        designed for individual use (desktops, laptops).
  \item \textbf{Servers} --- computers that provide services to other
        computers over a network; optimized for reliability, throughput, and
        availability.
  \item \textbf{Embedded computers} --- processors built into other devices
        (cars, appliances, IoT sensors); typically have strict power, cost,
        and real-time constraints.
\end{enumerate}

Other valid answers include supercomputers, warehouse-scale computers, and
personal mobile devices (PMDs).

%-----------------------------------------------------------------------------
\subsection*{1.2 --- Seven Great Ideas Matching}
\hypertarget{sol-1.2}{}
\markright{1.2}
\addcontentsline{toc}{subsection}{1.2 --- Seven Great Ideas Matching}

\begin{center}
\begin{tabular}{cl}
\toprule
\textbf{Item} & \textbf{Matching Idea} \\
\midrule
a. Assembly lines              & Performance via Pipelining \\
b. Suspension bridge cables    & Dependability via Redundancy \\
c. Navigation with wind info   & Performance via Prediction \\
d. Express elevators           & Make the Common Case Fast \\
e. Library reserve desk        & Hierarchy of Memories \\
f. Wider CMOS gate area        & Performance via Parallelism \\
g. Self-driving / existing sensors & Use Abstraction to Simplify Design \\
\bottomrule
\end{tabular}
\end{center}

\textbf{Reasoning:}
\begin{itemize}
  \item[a.] Each station performs one step, handing the product to the next
        --- classic pipelining.
  \item[b.] Many cables carry the load; if one fails, others compensate
        --- redundancy.
  \item[c.] The system predicts future position by incorporating wind
        --- prediction.
  \item[d.] Express elevators skip uncommon floors, making the most frequent
        trips fast --- common case fast.
  \item[e.] Popular books are kept closer (reserve desk) while others are in
        the stacks --- memory hierarchy.
  \item[f.] A wider transistor gate allows more current to flow in parallel
        --- parallelism.
  \item[g.] Reusing an existing sensor subsystem without redesigning it
        --- abstraction.
\end{itemize}

%-----------------------------------------------------------------------------
\subsection*{1.3 --- High-Level Language to Machine Code}
\hypertarget{sol-1.3}{}
\markright{1.3}
\addcontentsline{toc}{subsection}{1.3 --- High-Level Language to Machine Code}

\begin{enumerate}
  \item \textbf{Compiler} translates the high-level language program (e.g.\ C)
        into \textbf{assembly language}.
  \item \textbf{Assembler} translates the assembly language into
        \textbf{machine language} (object module, binary).
  \item \textbf{Linker} combines one or more object modules with library
        routines to produce the \textbf{executable}.
  \item \textbf{Loader} loads the executable into memory for the processor
        to fetch and execute.
\end{enumerate}

%-----------------------------------------------------------------------------
\subsection*{1.4 --- Color Display Frame Buffer}
\hypertarget{sol-1.4}{}
\markright{1.4}
\addcontentsline{toc}{subsection}{1.4 --- Color Display Frame Buffer}

$640 \times 480$ pixels, 8 bits/color $\times$ 3 colors = 24 bits/pixel = 3 bytes/pixel.

\medskip\noindent\textbf{a.} Frame buffer size:
\[
  640 \times 480 \times 3 = 921{,}600 \text{ bytes} \approx 900 \text{ KB}
\]

\noindent\textbf{b.} Time over 100 Mbit/s network:
\[
  921{,}600 \times 8 = 7{,}372{,}800 \text{ bits}
\]
\[
  t = \frac{7{,}372{,}800}{100 \times 10^6} \approx 0.0737 \text{ s}
  \approx 73.7 \text{ ms}
\]

%-----------------------------------------------------------------------------
\subsection*{1.5 --- Comparing Three Processors}
\hypertarget{sol-1.5}{}
\markright{1.5}
\addcontentsline{toc}{subsection}{1.5 --- Comparing Three Processors}

\noindent\textbf{a.} Instructions per second $= \text{Clock rate} / \text{CPI}$:
\[
  \text{P1} = \frac{3 \times 10^9}{1.5} = 2.0 \times 10^9, \quad
  \text{P2} = \frac{2.5 \times 10^9}{1.0} = 2.5 \times 10^9, \quad
  \text{P3} = \frac{4.0 \times 10^9}{2.2} \approx 1.82 \times 10^9
\]
\textbf{P2} has the highest performance.

\medskip\noindent\textbf{b.} Execute in 10 s $\Rightarrow$
Cycles $= \text{Clock rate} \times 10$, \quad
Instructions $= \text{Cycles} / \text{CPI}$:

\begin{center}
\begin{tabular}{lccc}
\toprule
& P1 & P2 & P3 \\
\midrule
Cycles ($\times 10^9$)       & 30   & 25    & 40 \\
Instructions ($\times 10^9$) & 20   & 25    & 18.18 \\
\bottomrule
\end{tabular}
\end{center}

\medskip\noindent\textbf{c.} New time $= 0.7 \times 10 = 7$ s, new CPI $= 1.2
\times \text{CPI}_{\text{old}}$.
From the original: $\text{IC} \times \text{CPI}_{\text{old}} =
10 \times \text{Clock}_{\text{old}}$.
\[
  \text{Clock}_{\text{new}}
  = \frac{\text{IC} \times 1.2 \times \text{CPI}_{\text{old}}}{7}
  = \frac{1.2 \times 10 \times \text{Clock}_{\text{old}}}{7}
  = \frac{12}{7}\;\text{Clock}_{\text{old}}
\]

\begin{center}
\begin{tabular}{lcc}
\toprule
& Old Clock (GHz) & Required Clock (GHz) \\
\midrule
P1 & 3.0 & $\frac{12}{7}\times 3.0 \approx 5.14$ \\
P2 & 2.5 & $\frac{12}{7}\times 2.5 \approx 4.29$ \\
P3 & 4.0 & $\frac{12}{7}\times 4.0 \approx 6.86$ \\
\bottomrule
\end{tabular}
\end{center}

%-----------------------------------------------------------------------------
\subsection*{1.6 --- Intel Processor Performance Trends}
\hypertarget{sol-1.6}{}
\markright{1.6}
\addcontentsline{toc}{subsection}{1.6 --- Intel Processor Performance Trends}

Average annual rate: $r = (V_n / V_0)^{1/n} - 1$ where $n = 9$ years (2010--2019).

Years to double: $n_d = \ln 2 / \ln(1 + r)$.

\begin{center}
\begin{tabular}{lcccc}
\toprule
\textbf{Indicator} & \textbf{2010} & \textbf{2019} & \textbf{Rate/yr} &
\textbf{Yrs to double} \\
\midrule
Tech (nm, smaller = better) & 45 & 14 & 13.8\% & 5.4 \\
Clock (GHz)      & 2.93 & 3.60 & 2.3\%  & 30.2 \\
Integer IPC/Core & 1.0  & 1.5  & 4.6\%  & 15.4 \\
Cores            & 4    & 8    & 8.0\%  & 9.0  \\
Max DRAM (GB)    & 16   & 128  & 26.0\% & 3.0  \\
DRAM BW (GB/s)   & 21   & 41.6 & 7.9\%  & 9.1  \\
\bottomrule
\end{tabular}
\end{center}

For Tech, ``improvement'' means shrinking: ratio $= 45/14 \approx 3.21$.

%-----------------------------------------------------------------------------
\subsection*{1.7 --- Two ISA Implementations Compared}
\hypertarget{sol-1.7}{}
\markright{1.7}
\addcontentsline{toc}{subsection}{1.7 --- Two ISA Implementations Compared}

Weighted CPI:
\begin{align*}
  \text{CPI}_{\text{P1}} &= 0.10(1) + 0.20(2) + 0.50(3) + 0.20(3) = 2.6 \\
  \text{CPI}_{\text{P2}} &= 0.10(2) + 0.20(2) + 0.50(2) + 0.20(2) = 2.0
\end{align*}

\noindent\textbf{a.} Execution time $= \text{IC} \times \text{CPI} /
\text{Clock}$ with $\text{IC} = 10^6$:
\begin{align*}
  T_{\text{P1}} &= \frac{10^6 \times 2.6}{2.5 \times 10^9}
                 = 1.04 \times 10^{-3}\text{ s} = 1.04\text{ ms} \\
  T_{\text{P2}} &= \frac{10^6 \times 2.0}{3.0 \times 10^9}
                 = 0.667 \times 10^{-3}\text{ s} = 0.667\text{ ms}
\end{align*}
\textbf{P2 is faster} by a factor of $1.04/0.667 \approx 1.56$.

\medskip\noindent\textbf{b.} Clock cycles:
\[
  \text{P1}: 10^6 \times 2.6 = 2.6 \times 10^6, \qquad
  \text{P2}: 10^6 \times 2.0 = 2.0 \times 10^6
\]

%-----------------------------------------------------------------------------
\subsection*{1.8 --- Compiler Impact on Performance}
\hypertarget{sol-1.8}{}
\markright{1.8}
\addcontentsline{toc}{subsection}{1.8 --- Compiler Impact on Performance}

Clock cycle time $= 1$ ns $\Rightarrow$ Clock rate $= 1$ GHz $= 10^9$ Hz.

\medskip\noindent\textbf{a.} $\text{CPU time} = \text{IC} \times \text{CPI}
\times \text{cycle time}$:
\begin{align*}
  \text{A}: &\quad 1.1 = 10^9 \times \text{CPI}_A \times 10^{-9}
    \;\Rightarrow\; \text{CPI}_A = 1.1 \\
  \text{B}: &\quad 1.5 = 1.2 \times 10^9 \times \text{CPI}_B \times 10^{-9}
    \;\Rightarrow\; \text{CPI}_B = 1.25
\end{align*}

\noindent\textbf{b.} Equal execution times on different processors:
\[
  \frac{\text{IC}_A \times \text{CPI}_A}{\text{Clock}_A}
  = \frac{\text{IC}_B \times \text{CPI}_B}{\text{Clock}_B}
  \quad\Rightarrow\quad
  \frac{\text{Clock}_A}{\text{Clock}_B}
  = \frac{10^9 \times 1.1}{1.2 \times 10^9 \times 1.25}
  = \frac{1.1}{1.5} \approx 0.733
\]
Compiler A's processor clock is about 73.3\% of compiler B's.

\medskip\noindent\textbf{c.} New compiler: IC $= 6 \times 10^8$, CPI $= 1.1$:
\[
  T_{\text{new}} = 6 \times 10^8 \times 1.1 \times 10^{-9} = 0.66\text{ s}
\]
\[
  \text{Speedup vs.\ A} = \frac{1.1}{0.66} \approx 1.67, \qquad
  \text{Speedup vs.\ B} = \frac{1.5}{0.66} \approx 2.27
\]

%-----------------------------------------------------------------------------
\subsection*{1.9 --- Processor Power and Capacitance}
\hypertarget{sol-1.9}{}
\markright{1.9}
\addcontentsline{toc}{subsection}{1.9 --- Processor Power and Capacitance}

Dynamic power: $P_{\text{dyn}} = C \times V^2 \times f$.

\phantomsection\addcontentsline{toc}{subsubsection}{1.9.1}
\medskip\noindent\textbf{1.9.1} Solve for $C$:

\begin{center}
\begin{tabular}{lcc}
\toprule
& Prescott & Ivy Bridge \\
\midrule
$P_{\text{dyn}}$ (W)  & 90   & 40 \\
$V$ (V)                & 1.25 & 0.9 \\
$f$ (GHz)              & 3.6  & 3.4 \\
\midrule
$C = P_{\text{dyn}}/(V^2 f)$ & $\dfrac{90}{1.5625 \times 3.6\text{G}}
  \approx 16.0$ nF & $\dfrac{40}{0.81 \times 3.4\text{G}} \approx 14.5$ nF \\
\bottomrule
\end{tabular}
\end{center}

\phantomsection\addcontentsline{toc}{subsubsection}{1.9.2}
\medskip\noindent\textbf{1.9.2} Static power fraction:

\begin{center}
\begin{tabular}{lcc}
\toprule
& Prescott & Ivy Bridge \\
\midrule
$P_{\text{total}}$ (W)       & 100  & 70 \\
$P_{\text{static}}$ (W)      & 10   & 30 \\
Static \%                     & 10\% & 42.9\% \\
$P_{\text{static}}/P_{\text{dyn}}$ & $1/9 \approx 0.111$ & $3/4 = 0.75$ \\
\bottomrule
\end{tabular}
\end{center}

\phantomsection\addcontentsline{toc}{subsubsection}{1.9.3}
\medskip\noindent\textbf{1.9.3}
Let $V' = V - \Delta V$.  Leakage current $I_{\text{leak}}$ is constant,
so $P_{\text{static}}' = V' \times I_{\text{leak}}$, and
$P_{\text{dyn}}' = C V'^{2} f$.

Set $P' = 0.9 \times P_{\text{total}}$ and solve the quadratic
$C f \cdot V'^2 + I_{\text{leak}} \cdot V' - 0.9 P_{\text{total}} = 0$.

\textbf{Prescott:} $I_{\text{leak}} = 10/1.25 = 8$ A, $Cf = 57.6$ A/V.
\[
  57.6\,V'^2 + 8\,V' - 90 = 0
  \;\Rightarrow\;
  V' = \frac{-8 + \sqrt{64 + 4(57.6)(90)}}{2(57.6)}
     = \frac{-8 + 144.2}{115.2} \approx 1.182\text{ V}
\]
Reduction $\approx 0.068$ V ($5.4\%$).

\textbf{Ivy Bridge:} $I_{\text{leak}} = 30/0.9 = 33.3$ A, $Cf = 49.4$ A/V.
\[
  49.4\,V'^2 + 33.3\,V' - 63 = 0
  \;\Rightarrow\;
  V' = \frac{-33.3 + \sqrt{33.3^2 + 4(49.4)(63)}}{2(49.4)}
     = \frac{-33.3 + 116.4}{98.8} \approx 0.841\text{ V}
\]
Reduction $\approx 0.059$ V ($6.5\%$).

%-----------------------------------------------------------------------------
\subsection*{1.10 --- Multiprocessor Execution Speedup}
\hypertarget{sol-1.10}{}
\markright{1.10}
\addcontentsline{toc}{subsection}{1.10 --- Multiprocessor Execution Speedup}

CPI: arithmetic = 1, load/store = 12, branch = 5.

Total clock cycles on a single processor:
\[
  2.56 \times 10^9 \times 1 + 1.28 \times 10^9 \times 12
  + 256 \times 10^6 \times 5
  = (2.56 + 15.36 + 1.28) \times 10^9
  = 19.2 \times 10^9
\]

Assuming a 2 GHz clock: $T_1 = 19.2 \times 10^9 / (2 \times 10^9) = 9.6$ s.

\phantomsection\addcontentsline{toc}{subsubsection}{1.10.1}
\medskip\noindent\textbf{1.10.1}
Assuming instructions are divided equally among $p$ processors:
\[
  T_p = \frac{T_1}{p} = \frac{9.6}{p}
\]

\begin{center}
\begin{tabular}{lccc}
\toprule
$p$ & $T_p$ (s) & Speedup \\
\midrule
1 & 9.6  & 1.0 \\
2 & 4.8  & 2.0 \\
8 & 1.2  & 8.0 \\
\bottomrule
\end{tabular}
\end{center}

(With perfectly parallelizable work, speedup is linear.)

\phantomsection\addcontentsline{toc}{subsubsection}{1.10.2}
\medskip\noindent\textbf{1.10.2}
Arithmetic time halved $\Rightarrow$ new arithmetic cycles
$= 2.56 \times 10^9 \times 0.5 = 1.28 \times 10^9$.
\[
  \text{New total} = 1.28 + 15.36 + 1.28 = 17.92 \times 10^9 \text{ cycles}
\]
\[
  \text{Speedup} = \frac{19.2}{17.92} \approx 1.071
\]
Arithmetic as percentage of new total:
$1.28 / 17.92 \approx 7.1\%$ (down from $2.56/19.2 = 13.3\%$).

%-----------------------------------------------------------------------------
\subsection*{1.11 --- Wafer Yield and Cost per Die}
\hypertarget{sol-1.11}{}
\markright{1.11}
\addcontentsline{toc}{subsection}{1.11 --- Wafer Yield and Cost per Die}

Die area $\approx$ Wafer area / Dies per wafer.

Wafer yield formula (from textbook):
$\text{Yield} = 1 / (1 + \text{Defects/cm}^2 \times \text{Die area}/2)^2$.

\begin{center}
\begin{tabular}{lcc}
\toprule
& Wafer 1 & Wafer 2 \\
\midrule
Diameter (cm) & 15 & 20 \\
Wafer area (cm$^2$) & $\pi(7.5)^2 \approx 176.7$ & $\pi(10)^2 \approx 314.2$ \\
Dies & 84 & 100 \\
Die area (cm$^2$) & $\approx 2.10$ & $\approx 3.14$ \\
Defects/cm$^2$ & 0.020 & 0.031 \\
\bottomrule
\end{tabular}
\end{center}

\phantomsection\addcontentsline{toc}{subsubsection}{1.11.1}
\medskip\noindent\textbf{1.11.1} Yield:
\begin{align*}
  Y_1 &= \frac{1}{(1 + 0.020 \times 2.10/2)^2}
       = \frac{1}{(1.021)^2} \approx 0.959 \quad (95.9\%) \\
  Y_2 &= \frac{1}{(1 + 0.031 \times 3.14/2)^2}
       = \frac{1}{(1.0487)^2} \approx 0.909 \quad (90.9\%)
\end{align*}

\noindent\textbf{1.11.2} Cost per die $=$ Wafer cost / (Dies $\times$ Yield):
\begin{align*}
  C_1 &= \$12 / (84 \times 0.959) = \$12 / 80.6 \approx \$0.149 \\
  C_2 &= \$15 / (100 \times 0.909) = \$15 / 90.9 \approx \$0.165
\end{align*}

\noindent\textbf{1.11.3}
Dies increase by 10\% $\Rightarrow$ new die area $=$ wafer area / new dies:
\begin{align*}
  A_1' &= 176.7 / (84 \times 1.1) = 176.7 / 92.4 \approx 1.91\text{ cm}^2 \\
  A_2' &= 314.2 / (100 \times 1.1) = 314.2 / 110 \approx 2.86\text{ cm}^2
\end{align*}

\noindent\textbf{1.11.4}
Solve for defect density $D$ from yield $= 1/(1 + D \cdot A/2)^2$:
\[
  1 + \frac{DA}{2} = \frac{1}{\sqrt{Y}}
  \quad\Rightarrow\quad
  D = \frac{2}{A}\!\left(\frac{1}{\sqrt{Y}} - 1\right)
\]
Using Wafer 1's die area ($A = 2.10$ cm$^2$) as example:
\begin{align*}
  Y = 0.92: &\quad D = \frac{2}{2.10}(1/\sqrt{0.92} - 1)
             = 0.952 \times 0.04256 \approx 0.0405 \text{ defects/cm}^2 \\
  Y = 0.95: &\quad D = \frac{2}{2.10}(1/\sqrt{0.95} - 1)
             = 0.952 \times 0.02598 \approx 0.0247 \text{ defects/cm}^2
\end{align*}

%-----------------------------------------------------------------------------
\subsection*{1.12 --- SPEC CPU2006 Benchmark Analysis}
\hypertarget{sol-1.12}{}
\markright{1.12}
\addcontentsline{toc}{subsection}{1.12 --- SPEC CPU2006 Benchmark Analysis}

IC $= 2.389 \times 10^{12}$, Execution time $= 750$ s, Reference time
$= 9650$ s.

\phantomsection\addcontentsline{toc}{subsubsection}{1.12.1}
\medskip\noindent\textbf{1.12.1}
Clock cycle time $= 0.333$ ns $\Rightarrow f = 1/0.333\text{n} = 3.0$ GHz.
\[
  \text{CPI} = \frac{T \times f}{\text{IC}}
  = \frac{750 \times 3.0 \times 10^9}{2.389 \times 10^{12}}
  = \frac{2250 \times 10^9}{2389 \times 10^9} \approx 0.942
\]

\noindent\textbf{1.12.2}
\[
  \text{SPECratio} = \frac{\text{Reference time}}{\text{Execution time}}
  = \frac{9650}{750} \approx 12.87
\]

\noindent\textbf{1.12.3}
New clock $= 4$ GHz, IC reduced 15\% $\Rightarrow \text{IC}' = 0.85 \times
2.389 \times 10^{12} = 2.031 \times 10^{12}$.
SPECratio $= 13.7 \Rightarrow T' = 9650/13.7 \approx 704.4$ s.
\[
  \text{CPI}' = \frac{704.4 \times 4 \times 10^9}{2.031 \times 10^{12}}
  = \frac{2817.6}{2031} \approx 1.387
\]

\noindent\textbf{1.12.4}
IC increases 10\%, CPI and clock unchanged:
$T' = 1.1 \times T = 1.1 \times 750 = 825$ s.
\textbf{Increase = 10\%.}

\phantomsection\addcontentsline{toc}{subsubsection}{1.12.5}
\medskip\noindent\textbf{1.12.5}
IC $+10\%$, CPI $-15\%$, clock unchanged:
\[
  T' = \frac{1.1 \times \text{IC} \times 0.85 \times \text{CPI}}{f}
     = 0.935 \times T = 0.935 \times 750 = 701.3\text{ s}
\]
CPU time \textbf{decreases by 6.5\%}.

\phantomsection\addcontentsline{toc}{subsubsection}{1.12.6}
\medskip\noindent\textbf{1.12.6}
IC $+10\%$, CPI $-15\%$, and require $T' = 0.8T$ (20\% reduction).
\[
  T' = \frac{\text{IC}' \times \text{CPI}'}{f'}
  \;\Rightarrow\;
  f' = \frac{1.1 \times 0.85 \times \text{IC} \times \text{CPI}}{0.8\,T}
     = \frac{0.935}{0.8}\,f = 1.169\,f
\]
Clock must \textbf{increase by $\approx 16.9\%$}
(e.g.\ $3.0 \to 3.51$ GHz).

%-----------------------------------------------------------------------------
\subsection*{1.13 --- Performance Equation Pitfall}
\hypertarget{sol-1.13}{}
\markright{1.13}
\addcontentsline{toc}{subsection}{1.13 --- Performance Equation Pitfall}

P1: $f = 4$ GHz, CPI $= 0.9$, IC $= 5.0 \times 10^9$. \quad
P2: $f = 3$ GHz, CPI $= 0.75$, IC $= 1.0 \times 10^9$.

\phantomsection\addcontentsline{toc}{subsubsection}{1.13.1}
\medskip\noindent\textbf{1.13.1}
\begin{align*}
  T_{\text{P1}} &= \frac{5.0 \times 10^9 \times 0.9}{4 \times 10^9}
                 = 1.125\text{ s} \\
  T_{\text{P2}} &= \frac{1.0 \times 10^9 \times 0.75}{3 \times 10^9}
                 = 0.25\text{ s}
\end{align*}
\textbf{P2 is faster} despite having a lower clock rate.  The fallacy is
\textbf{false}.

\phantomsection\addcontentsline{toc}{subsubsection}{1.13.2}
\medskip\noindent\textbf{1.13.2}
P1 time for $10^9$ instructions:
$T = 10^9 \times 0.9 / (4 \times 10^9) = 0.225$ s.

P2 instructions in same time:
$0.225 = \text{IC}_{P2} \times 0.75 / (3 \times 10^9)$
$\Rightarrow \text{IC}_{P2} = 0.225 \times 3 \times 10^9 / 0.75
= 9.0 \times 10^8$.

P2 executes \textbf{fewer} instructions ($9.0 \times 10^8 < 1.0 \times 10^9$)
in the same time.  More instructions $\neq$ more time.

\phantomsection\addcontentsline{toc}{subsubsection}{1.13.3}
\medskip\noindent\textbf{1.13.3}
$\text{MIPS} = f / (\text{CPI} \times 10^6)$:
\[
  \text{P1} = \frac{4 \times 10^9}{0.9 \times 10^6} = 4444, \qquad
  \text{P2} = \frac{3 \times 10^9}{0.75 \times 10^6} = 4000
\]
P1 has higher MIPS yet P2 is faster (from 1.13.1).  MIPS ignores differences
in instruction count and is therefore a \textbf{misleading} metric for
cross-program or cross-ISA comparisons.

%-----------------------------------------------------------------------------
\subsection*{1.14 --- Amdahl's Law and Partial Improvement}
\hypertarget{sol-1.14}{}
\markright{1.14}
\addcontentsline{toc}{subsection}{1.14 --- Amdahl's Law and Partial Improvement}

Total $= 250$ s.  FP $= 70$ s, L/S $= 85$ s, Branch $= 40$ s.
INT $= 250 - 70 - 85 - 40 = 55$ s.

\phantomsection\addcontentsline{toc}{subsubsection}{1.14.1}
\medskip\noindent\textbf{1.14.1}
FP reduced 20\%: new FP $= 70 \times 0.8 = 56$ s.
\[
  T' = 56 + 85 + 40 + 55 = 236\text{ s} \quad
  \text{(reduction of } 14\text{ s} = 5.6\%)
\]

\noindent\textbf{1.14.2}
Need $T' = 0.8 \times 250 = 200$ s $\Rightarrow$ must save 50 s.
FP, L/S, Branch are unchanged, so all 50 s must come from INT:
new INT $= 55 - 50 = 5$ s.
INT must be reduced by $50/55 \approx$ \textbf{90.9\%}.

\phantomsection\addcontentsline{toc}{subsubsection}{1.14.3}
\medskip\noindent\textbf{1.14.3}
Need 50 s reduction from branch alone, but branch $= 40$ s total.
Even reducing branch to 0 s saves only 40 s $< 50$ s.
\textbf{No}, it is impossible.

%-----------------------------------------------------------------------------
\subsection*{1.15 --- CPI Optimization for Speedup}
\hypertarget{sol-1.15}{}
\markright{1.15}
\addcontentsline{toc}{subsection}{1.15 --- CPI Optimization for Speedup}

\begin{center}
\begin{tabular}{lcccc}
\toprule
Type & IC ($\times 10^6$) & CPI & Cycles ($\times 10^6$) \\
\midrule
FP     & 50  & 1 & 50  \\
INT    & 110 & 1 & 110 \\
L/S    & 80  & 4 & 320 \\
Branch & 16  & 2 & 32  \\
\midrule
Total  &     &   & 512 \\
\bottomrule
\end{tabular}
\end{center}

$T = 512 \times 10^6 / (2 \times 10^9) = 0.256$ s.  Target for $2\times$
speedup: $256 \times 10^6$ cycles.

\phantomsection\addcontentsline{toc}{subsubsection}{1.15.1}
\medskip\noindent\textbf{1.15.1}
FP cycles needed $= 256 - 110 - 320 - 32 = -206 \times 10^6$.
This is negative --- \textbf{impossible} to achieve $2\times$ speedup by
improving FP CPI alone.

\phantomsection\addcontentsline{toc}{subsubsection}{1.15.2}
\medskip\noindent\textbf{1.15.2}
L/S cycles needed $= 256 - 50 - 110 - 32 = 64 \times 10^6$.
\[
  \text{New CPI}_{L/S} = \frac{64 \times 10^6}{80 \times 10^6} = 0.8
\]
L/S CPI must drop from 4 to 0.8 --- an \textbf{80\% reduction}.

\phantomsection\addcontentsline{toc}{subsubsection}{1.15.3}
\medskip\noindent\textbf{1.15.3}
FP and INT CPI $\times 0.6$; L/S and Branch CPI $\times 0.7$:
\begin{align*}
  \text{FP}' &= 50 \times 10^6 \times (1 \times 0.6) = 30 \times 10^6 \\
  \text{INT}' &= 110 \times 10^6 \times (1 \times 0.6) = 66 \times 10^6 \\
  \text{L/S}' &= 80 \times 10^6 \times (4 \times 0.7) = 224 \times 10^6 \\
  \text{Branch}' &= 16 \times 10^6 \times (2 \times 0.7) = 22.4 \times 10^6
\end{align*}
\[
  \text{Total}' = 342.4 \times 10^6, \qquad
  T' = 342.4 / 2000 = 0.1712\text{ s}
\]
\[
  \text{Speedup} = \frac{0.256}{0.1712} \approx 1.50\times
\]

%-----------------------------------------------------------------------------
\subsection*{1.16 --- Parallel Overhead and Speedup}
\hypertarget{sol-1.16}{}
\markright{1.16}
\addcontentsline{toc}{subsection}{1.16 --- Parallel Overhead and Speedup}

$t = 100$ s single-processor, overhead $= 4$ s/processor.
Per-processor time $= 100/p + 4$.

\[
  \text{Speedup} = \frac{100}{100/p + 4}, \qquad
  \text{Ideal speedup} = p, \qquad
  \text{Ratio} = \frac{\text{Speedup}}{p}
\]

\begin{center}
\begin{tabular}{rrrr}
\toprule
$p$ & Time (s) & Speedup & Ratio \\
\midrule
2   & 54.000 & 1.85  & 0.926 \\
4   & 29.000 & 3.45  & 0.862 \\
8   & 16.500 & 6.06  & 0.758 \\
16  & 10.250 & 9.76  & 0.610 \\
32  & 7.125  & 14.04 & 0.439 \\
64  & 5.563  & 17.98 & 0.281 \\
128 & 4.781  & 20.92 & 0.163 \\
\bottomrule
\end{tabular}
\end{center}

As $p$ grows, overhead dominates and speedup saturates.
The ratio shows that efficiency drops sharply beyond $\sim$16 processors
for this workload.
